```latex
\documentclass[11pt]{article}
\usepackage{amsmath,amssymb,amsthm}
\usepackage[margin=1.2in]{geometry}

\newtheorem{theorem}{Theorem}[section]
\newtheorem{lemma}[theorem]{Lemma}
\newtheorem{proposition}[theorem]{Proposition}
\newtheorem{corollary}[theorem]{Corollary}
\theoremstyle{definition}
\newtheorem{definition}[theorem]{Definition}
\newtheorem{convention}[theorem]{Convention}
\newtheorem{example}[theorem]{Example}
\theoremstyle{remark}
\newtheorem{remark}[theorem]{Remark}

\numberwithin{equation}{section}

\begin{document}

\title{Surgeries on $L$-space knots have at most two prime summands}
\author{}
\date{}
\maketitle

\begin{abstract}
We prove that if $K\subset S^{3}$ is an $L$-space knot, then every closed,
oriented $3$-manifold obtained by Dehn surgery on $K$ has at most two
prime summands.  More precisely, such a surgery is either irreducible,
or else $K$ is a cable $C_{p,q}(J)$, the surgery coefficient is the cabling
slope $pq$, and the surgered manifold is the connected sum
\[
S^{3}_{pq}\bigl(C_{p,q}(J)\bigr)\;\cong\;L(p,q)\,\#\,S^{3}_{q/p}(J),
\]
whose two summands are prime: the first is a lens space, and the second is
irreducible by the classical theorem of Culler--Gordon--Luecke--Shalen that
non-integral surgeries on nontrivial knots are irreducible.  The key
topological input is the theorem of Hedden--Levine that the cabling
conjecture holds for $L$-space knots.  The answer to the question in the
title is therefore \emph{yes}.
\end{abstract}

\section{Introduction}\label{sec:intro}

Recall that a knot $K\subset S^{3}$ is an \emph{$L$-space knot} if some
Dehn surgery on $K$ produces an $L$-space in the sense of Ozsv\'ath and
Szab\'o's Heegaard Floer homology (Definition~\ref{def:lspace} below).
The class of $L$-space knots contains all torus knots and, more generally,
all sufficiently large cables of $L$-space knots, and it also contains
infinitely many hyperbolic knots.

For a general knot $K\subset S^{3}$, the structure of \emph{reducible}
surgeries---surgeries producing a disconnected sum $Y_{1}\#Y_{2}$ with
both summands nontrivial---is governed by the classical \emph{cabling
conjecture}, which predicts that a reducible surgery occurs only when $K$
is a cable knot and the surgery coefficient is the cabling slope, in which
case the output has exactly two prime summands, one of which is a lens
space.  The cabling conjecture remains open in general.  However, for
$L$-space knots it is a theorem of Hedden and Levine
\cite{HL16}, obtained from their bordered Heegaard Floer analysis of
spliced manifolds.  Combining this theorem with the classical Dehn surgery
calculus for cables and with the foundational results of
Culler--Gordon--Luecke--Shalen on non-integral surgeries, we obtain a
complete answer to the question in the title.

\begin{theorem}[Main Theorem]\label{thm:main-intro}
Let $K\subset S^{3}$ be an $L$-space knot, and let $r\in\mathbb{Q}\cup\{\infty\}$
be a surgery slope.  Then the closed, oriented $3$-manifold
$Y=S^{3}_{r}(K)$ has at most two prime summands.  More precisely:
\begin{enumerate}
\item[(i)] if $Y$ is irreducible, then $Y$ is prime, so $Y$ has one prime
summand (or zero, if $Y\cong S^{3}$);
\item[(ii)] if $Y$ is reducible, then $K$ is a cable knot: there is a knot
$J\subset S^{3}$ and coprime integers $p\ge 2$, $q\neq 0$ with
$K=C_{p,q}(J)$ and $r=pq$, and
\[
S^{3}_{pq}\bigl(C_{p,q}(J)\bigr)\;\cong\;L(p,q)\;\#\;S^{3}_{q/p}(J),
\]
where $L(p,q):=S^{3}_{p/q}(O)$ is a lens space and the second summand
$S^{3}_{q/p}(J)$ is irreducible.  In particular $Y$ has exactly two prime
summands.
\end{enumerate}
\end{theorem}

Thus the answer to the question posed above is \textbf{yes}.

\paragraph{Plan of the paper.}
Section~\ref{sec:prelim} fixes conventions for Dehn surgery, recalls the
definitions of prime, reducible, and irreducible $3$-manifolds together
with the prime decomposition theorem, and records the basic examples
(surgeries on the unknot) and elementary lemmas that are used later,
including the (proved) fact that irreducible implies prime.
Section~\ref{sec:cabling-conjecture} states the Hedden--Levine theorem
(the cabling conjecture for $L$-space knots) and describes its mechanism.
Section~\ref{sec:cables} records the classical surgery calculus for cable
knots: Moser's theorem for torus knots and the cabling-slope splitting
formula for general cables, with a verification of the slope bookkeeping.
Section~\ref{sec:proof} assembles the proof of
Theorem~\ref{thm:main-intro}.  Section~\ref{sec:remarks} contains
corollaries, examples, and remarks, including a precise statement of the
``exactly two summands'' refinement.

\section{Preliminaries and conventions}\label{sec:prelim}

\subsection{Dehn surgery}

Throughout, all manifolds are compact, oriented, and connected unless
stated otherwise, and all knots are oriented.

\begin{convention}[Slopes and surgery]\label{conv:surgery}
Let $K\subset S^{3}$ be a knot with exterior
$E_{K}=S^{3}\setminus\operatorname{int}N(K)$.  A \emph{slope} on
$\partial E_{K}$ is a primitive class in $H_{1}(\partial E_{K})\cong
\mathbb{Z}^{2}$.  Fix a meridian $\mu_{K}$ and the Seifert longitude
$\lambda_{K}$ (the null-homologous class with
$\operatorname{lk}(\lambda_{K},K)=0$), with
$\mu_{K}\cdot\lambda_{K}=1$.  For relatively prime integers $p,q$ with
$q>0$, the \emph{$p/q$-slope} is the class
\[
q\,[\lambda_{K}]+p\,[\mu_{K}]\in H_{1}(\partial E_{K}),
\]
and the \emph{$p/q$-Dehn surgery} is
\[
S^{3}_{p/q}(K):=E_{K}\cup\bigl(S^{1}\times D^{2}\bigr),
\]
glued so that the meridian $\{\text{pt}\}\times\partial D^{2}$ of the
attached solid torus is identified with a curve of $p/q$-slope.  We write
$S^{3}_{n}(K):=S^{3}_{n/1}(K)$ for $n\in\mathbb{Z}$, and
$S^{3}_{\infty}(K):=S^{3}$.  A slope $p/q$ (in lowest terms) is
\emph{integral} if $q=1$ and \emph{non-integral} if $q\ge 2$.
\end{convention}

With this convention, the surgery curve of slope $r$ has linking number
$r$ with $K$, which is the classical meaning of the surgery coefficient.

\begin{lemma}\label{lem:H1}
For any knot $K\subset S^{3}$ and any slope $p/q$ in lowest terms,
\[
H_{1}\bigl(S^{3}_{p/q}(K)\bigr)\cong
\begin{cases}
\mathbb{Z}_{|p|}, & p\neq 0,\\[2pt]
\mathbb{Z}, & p=0.
\end{cases}
\]
Moreover $H_{1}(A\#B)\cong H_{1}(A)\oplus H_{1}(B)$ for closed, oriented
$3$-manifolds $A,B$.
\end{lemma}

\begin{proof}
Since $\lambda_{K}$ is null-homologous in $E_{K}$, we have
$H_{1}(E_{K})=\mathbb{Z}\langle\mu_{K}\rangle$ with
$[\lambda_{K}]=0$.  Mayer--Vietoris for the decomposition
$S^{3}_{p/q}(K)=E_{K}\cup(S^{1}\times D^{2})$ shows that
$H_{1}(S^{3}_{p/q}(K))$ is obtained from $H_{1}(E_{K})$ by killing the
image of the class of the surgery curve, which is
$p[\mu_{K}]$.  This gives $\mathbb{Z}_{|p|}$ for $p\neq0$ and
$\mathbb{Z}$ for $p=0$.  The Künneth/Mayer--Vietoris computation for
connected sums is standard and left to the reader.
\end{proof}

\subsection{Prime decompositions}

\begin{definition}\label{def:prime}
A closed, oriented $3$-manifold $Y$ is \emph{reducible} if it admits a
connected sum decomposition $Y\cong Y_{1}\#Y_{2}$ with
$Y_{i}\not\cong S^{3}$; equivalently (for $Y\not\cong S^{1}\times S^{2}$),
if $Y$ contains an embedded $2$-sphere that does not bound a $3$-ball.
$Y$ is \emph{irreducible} if every embedded $2$-sphere in $Y$ bounds a
$3$-ball.  $Y$ is \emph{prime} if $Y\cong Y_{1}\#Y_{2}$ implies
$Y_{1}\cong S^{3}$ or $Y_{2}\cong S^{3}$.
\end{definition}

By the prime decomposition theorem of Milnor and Kneser
\cite{Mi62,He76}, every closed, oriented, connected $3$-manifold $Y$ can
be written as
\[
Y\cong Y_{1}\#\cdots\# Y_{n}
\]
with each $Y_{i}$ prime and $Y_{i}\not\cong S^{3}$, and the summands
$Y_{1},\dots,Y_{n}$ are unique up to order and (after Perelman) up to
homeomorphism.  We say that \emph{$Y$ has $n$ prime summands}; thus
$S^{3}$ has zero prime summands.  In particular, the phrase ``$Y$ has at
most two prime summands'' is well defined.

\begin{lemma}\label{lem:irr-prime}
Let $Y$ be a closed, oriented, connected $3$-manifold.
\begin{enumerate}
\item[(i)] If $Y$ is irreducible, then $Y$ is prime.
\item[(ii)] $S^{1}\times S^{2}$ is reducible but prime.
\end{enumerate}
\end{lemma}

\begin{proof}
(i) Suppose $Y\cong Y_{1}\#Y_{2}$ and let $S$ be the separating
$2$-sphere of the connected sum.  If $S$ bounded a $3$-ball $B\subset Y$,
then one side of $S$ in $Y$ equals $B$; that side is, say,
$Y_{1}\setminus\operatorname{int}B_{1}$ after choosing the balls defining
the connected sum, so
$Y_{1}=(Y_{1}\setminus\operatorname{int}B_{1})\cup B_{1}\cong B\cup
B_{1}\cong S^{3}$, since gluing two $3$-balls along their boundary is
$S^{3}$.  Hence, if both $Y_{i}\not\cong S^{3}$, the sphere $S$ does not
bound a ball, contradicting irreducibility.
(ii) is standard; see \cite{He76}.
\end{proof}

\begin{lemma}\label{lem:lens}
Let $p,q$ be coprime with $p\neq 0$, and let
$L(p,q)$ denote the lens space $S^{3}_{p/q}(O)$, where $O\subset S^{3}$
is the unknot.  Then $L(p,q)$ is irreducible, hence prime.  More
generally, for every slope $r$, the manifold $S^{3}_{r}(O)$ is one of
$S^{3}$, $S^{1}\times S^{2}$, or a lens space; in all cases it is prime
or $S^{3}$.
\end{lemma}

\begin{proof}
Surgeries on the unknot are the classical lens spaces: an explicit
description of the exterior of $O$ as a solid torus shows that
$S^{3}_{p/q}(O)\cong L(p,q)$ for $p\neq 0$ and
$S^{3}_{0}(O)\cong S^{1}\times S^{2}$; see \cite[Chapter~9]{Ro90}.  Lens
spaces are irreducible \cite{He76,Ro90}.  Apply
Lemma~\ref{lem:irr-prime} and note that $S^{1}\times S^{2}$ is prime.
\end{proof}

\subsection{$L$-spaces and $L$-space knots}

\begin{definition}\label{def:lspace}
Let $Y$ be a closed, oriented $3$-manifold with
$b_{1}(Y)=0$.  Then $Y$ is an \emph{$L$-space} if
\[
\operatorname{rank}\widehat{HF}(Y)\;=\;\left|H_{1}(Y)\right|+1,
\]
i.e., if $\widehat{HF}(Y,\mathfrak{s})\cong\mathbb{Z}$ for each
$\operatorname{Spin}^{c}$-structure $\mathfrak{s}$; see
\cite{OSz04a,OSz04b}.  A knot $K\subset S^{3}$ is an \emph{$L$-space
knot} if $S^{3}_{r}(K)$ is an $L$-space for some slope $r$.  (The unknot
$O$ is an $L$-space knot under this definition; it will be treated
separately and is harmless.)
\end{definition}

We record for later use the fundamental ``large surgery'' property, due
to Ozsv\'ath and Szab\'o: if $K$ is a nontrivial $L$-space knot of genus
$g(K)$, then $S^{3}_{r}(K)$ is an $L$-space for every rational slope
$r\ge 2g(K)-1$ (and similarly for $r\le -(2g(K)-1)$); see
\cite{OSz05}.  Thus an $L$-space knot admits \emph{infinitely many}
$L$-space surgeries---a fact that is essential to the proof mechanism of
Theorem~\ref{thm:HL} below.

\subsection{Cables}\label{sec:cable-def}

\begin{convention}[Cables]\label{conv:cable}
Let $J\subset S^{3}$ be an oriented knot with meridian $\mu_{J}$ and
Seifert longitude $\lambda_{J}$, and let $T=\partial N(J)$.  For
integers $p\ge 2$ and $q\neq 0$ with $\gcd(p,|q|)=1$, the
\emph{$(p,q)$-cable} of $J$ is the oriented knot
\[
K=C_{p,q}(J)\subset T
\]
representing the class
$p[\lambda_{J}]+q[\mu_{J}]\in H_{1}(T)$, endowed with the orientation
induced from $J$.  The integer $p$ is the \emph{winding number}, the
integer $q$ the \emph{twisting number}; the \emph{cabling slope} of
$C_{p,q}(J)$ is $pq$ (with respect to the Seifert framing of the cable).
The \emph{cabling torus} is $T=\partial N(J)$, viewed as a torus in
$E_{K}$ after pushing $K$ slightly into $\operatorname{int}N(J)$.
\end{convention}

\begin{lemma}\label{lem:cable-basics}
With Convention~\ref{conv:cable}:
\begin{enumerate}
\item[(i)] the winding number is $p$: the class of $K$ in
$H_{1}(N(J))\cong\mathbb{Z}\langle[J]\rangle$ is $p[J]$, and
$\operatorname{lk}(K,J)=q$;
\item[(ii) ] $C_{p,q}(O)=T(p,q)$, the $(p,q)$-torus knot; in particular
$C_{p,q}(O)$ is nontrivial if and only if $p\ge2$ and $|q|\ge 2$;
\item[(iii)] the push-off of $K$ along the torus $T$ has linking number
$+pq$ with $K$; consequently the torus-framed longitude of $K$ equals
$\lambda_{K}+pq\,\mu_{K}$, which is exactly the $pq$-surgery curve of
Convention~\ref{conv:surgery};
\item[(iv)] the genus of $C_{p,q}(J)$ is
$g(C_{p,q}(J))=p\,g(J)+\tfrac{(p-1)(q-1)}{2}$
(Schubert's satellite genus formula).
\end{enumerate}
\end{lemma}

\begin{proof}
(i) In the solid torus $N(J)$ we have $H_{1}(N(J))=\mathbb{Z}\langle
\mu_{J}\text{-disk core}\rangle\cong\mathbb{Z}\langle[J]\rangle$,
$[\mu_{J}]=0$ (it bounds the meridian disk), and
$[\lambda_{J}]=[J]$ (it is a parallel copy of the core).  Hence
$[K]=p[\lambda_{J}]+q[\mu_{J}]=p[J]$.  Since
$\operatorname{lk}(\lambda_{J},J)=0$ and
$\operatorname{lk}(\mu_{J},J)=+1$, bilinearity of the linking number
gives $\operatorname{lk}(K,J)=q$.
(ii) is immediate from the definition: on $\partial N(O)$ the class
$p[\lambda_{O}]+q[\mu_{O}]$ is the standard $(p,q)$-torus curve.  The
nontriviality statement follows since $T(p,\pm1)$ and $T(\pm1,q)$ are
unknots.
(iii) Let $K^{+}$ be the push-off of $K$ along $T$.  The annulus
$A\subset T$ with $\partial A=K\cup K^{+}$ has core a curve of class
$p[\lambda_{J}]+q[\mu_{J}]$; the classical torus computation gives
$\operatorname{lk}(K,K^{+})=pq$ (the self-linking of the torus framing);
see \cite[Chapter~7]{Ro90}.  Since the Seifert longitude $\lambda_{K}$ is
by definition the unique class with $\operatorname{lk}(\lambda_{K},K)=0$
and $\lambda_{K}\cdot\mu_{K}=1$, we get
$K^{+}=\lambda_{K}+pq\,\mu_{K}$, which by
Convention~\ref{conv:surgery} is the surgery curve of slope $pq$.
(iv) is the classical satellite genus formula of Schubert; see
\cite{Ro90}, and \cite{EN85} for the equivalent splicing statement.  (We
include it only for orientation; it is not used in the proof of the
Main Theorem.)
\end{proof}

\section{The cabling conjecture for $L$-space knots}\label{sec:cabling-conjecture}

The following theorem is the Heegaard Floer-theoretic pillar of this
paper.  It is often summarized by saying that the cabling conjecture
holds for $L$-space knots.

\begin{theorem}[Hedden--Levine]\label{thm:HL}
Let $K\subset S^{3}$ be a nontrivial $L$-space knot, and suppose that
$S^{3}_{n}(K)$ is reducible for some slope $n\in\mathbb{Q}$.  Then $K$ is
a cable knot and $n$ is the cabling slope: there exist a knot
$J\subset S^{3}$ and coprime integers $p\ge2$, $q\neq0$ such that
\[
K=C_{p,q}(J)
\qquad\text{and}\qquad
n=pq .
\]
\end{theorem}

\begin{proof}[Discussion and attribution]
Theorem~\ref{thm:HL} is due to Hedden and Levine \cite{HL16}; the case of
torus knots (i.e., $J=O$) is classical and is due to Moser
\cite{Mos71}.  For the convenience of the reader we describe the
mechanism of the proof in \cite{HL16}, referring there for complete
details; the result is used below as a black box.

Suppose $S^{3}_{n}(K)\cong Y_{1}\#Y_{2}$ with
$Y_{i}\not\cong S^{3}$, and let $K^{*}$ denote the dual knot.  After an
isotopy, the separating sphere meets the surgery solid torus in
meridian disks and meets $K^{*}$ in two points, so that
$K^{*}=K^{*}_{1}\#K^{*}_{2}$ with $K^{*}_{i}\subset Y_{i}$.  Consequently
the exterior $E_{K}\cong E_{K^{*}}$ of the original knot decomposes along
the swallow-follow torus of the connected sum $K^{*}_{1}\#K^{*}_{2}$ into
two pieces, each a knot exterior with a $1$-handle attached---that is,
$E_{K}$ is a \emph{generalized splice} of the two exteriors, with
explicit winding data; this reduction is classical and is carried out in
the setting of splicing calculus in \cite{EN85}.  Now use that $K$ is an
$L$-space knot: by the large-surgery property recalled after
Definition~\ref{def:lspace}, infinitely many Dehn fillings of $E_{K}$
(namely all slopes $\ge 2g(K)-1$) are $L$-spaces.  The bordered Heegaard
Floer classification of spliced exteriors admitting infinitely many
$L$-space fillings, established in \cite{HL16}, then forces one of the
splice factors to be the exterior of a torus knot, with winding data
$(p,q)$ as in Convention~\ref{conv:cable}.  Unwinding the definitions
exhibits $K$ as the $(p,q)$-cable $C_{p,q}(J)$ of the knot $J$ carried by
the other factor, and forces the reducible slope to be the cabling slope
$n=pq$.
\end{proof}

\begin{remark}
Theorem~\ref{thm:HL} subsumes, in particular, the classical integrality
statement for reducible surgeries: the cabling slope $pq$ is an integer.
For general (non-$L$-space) knots the cabling conjecture remains open; the
proof of our Main Theorem therefore genuinely uses the Heegaard Floer
input.
\end{remark}

\section{The classical surgery calculus for cables}\label{sec:cables}

The second pillar is the classical description of surgery at a cabling
slope.  We first recall Moser's theorem.

\begin{theorem}[Moser {\cite{Mos71}}]\label{thm:moser}
Let $T(p,q)\subset S^{3}$ be the $(p,q)$-torus knot, $p,q$ coprime,
$p\ge2$, $|q|\ge 2$.  Then
\[
S^{3}_{pq}\bigl(T(p,q)\bigr)\;\cong\;L(p,q)\;\#\;L(q,p),
\]
and this is the unique reducible surgery on $T(p,q)$: for $n\neq pq$,
$S^{3}_{n}(T(p,q))$ is irreducible.
\end{theorem}

The general cable case is the following proposition.  We state it in a
symmetric form that also displays the answer to the slope bookkeeping:
the companion summand is filled along the class
$p[\lambda_{J}]+q[\mu_{J}]$---exactly the cabling class of
Convention~\ref{conv:cable}---which by Convention~\ref{conv:surgery} is
the surgery coefficient $q/p$.

\begin{proposition}[Cable splitting formula]\label{prop:cable-split}
Let $J\subset S^{3}$ be any knot (possibly the unknot), let $p\ge2$ and
$q\neq0$ be coprime, and let $K=C_{p,q}(J)$.  Then the cabling-slope
surgery satisfies
\[
S^{3}_{pq}(K)\;\cong\;S^{3}_{p/q}(O)\;\#\;S^{3}_{q/p}(J)
\;=\;L(p,q)\;\#\;S^{3}_{q/p}(J).
\]
Moreover, the cabling torus $T=\partial N(J)$ is essential in
$S^{3}_{pq}(K)$, and the two summands are the Dehn fillings of the two
pieces of $S^{3}_{pq}(K)\setminus T$.
\end{proposition}

\begin{proof}[Proof, following the classical calculus]
The statement is the standard cut-and-paste description of surgery at a
cabling slope; it is carried out in the Dehn surgery calculus for
satellite knots (see \cite[Chapter~9]{Ro90} and the splicing calculus of
\cite{EN85}), and it specializes to Theorem~\ref{thm:moser} when
$J=O$.  For the reader's convenience we recall the construction and then
verify the slope bookkeeping in detail, since the precise coefficient
$q/p$ is essential for what follows.

Push $K$ slightly into $\operatorname{int}N(J)$, so that the cabling
torus $T=\partial N(J)$ lies in $E_{K}$ and
\[
E_{K}\;=\;E_{J}\;\cup_{T}\;\mathcal{C},
\qquad
\mathcal{C}:=N(J)\setminus\operatorname{int}N(K)
\]
where $\mathcal{C}$ is the \emph{cable space}: the complement of the
$(p,q)$-curve in the solid torus $N(J)$, with boundary components $T$ and
$\partial N(K)$.  Then
\[
S^{3}_{pq}(K)=E_{J}\cup_{T}\mathcal{C}\cup Z,
\]
where $Z\cong S^{1}\times D^{2}$ is the surgery solid torus attached to
$\partial N(K)$ along the $pq$-slope.  By
Lemma~\ref{lem:cable-basics}(iii), the $pq$-surgery curve on
$\partial N(K)$ is the torus-framed longitude $\lambda^{T}$, i.e.\ the
push-off of $K$ along $T$.  Consequently the annulus
$B\subset\mathcal{C}$ obtained as the closure of
$T\setminus(\text{annular neighborhood of }K\text{ in }T)$ has its two
boundary components on $\partial N(K)$, each parallel to $\lambda^{T}$;
each bounds a meridian disk of $Z$.  Thus
\[
S:=B\cup D_{1}\cup D_{2}
\]
is an embedded $2$-sphere in $S^{3}_{pq}(K)$.  The classical analysis of
this sphere (cf.\ the proof of Theorem~\ref{thm:moser} in
\cite{Mos71}, and the satellite surgery calculus in
\cite{Ro90,EN85}) shows that $S$ is essential---indeed
$S^{3}_{pq}(K)\setminus S$ has two components, one of which is
$\bigl(\mathcal{C}\cup Z\bigr)$ capped off, and the other is $E_{J}$
capped off---and that capping the components yields
\[
S^{3}_{pq}(K)\;\cong\;\Bigl(\mathcal{C}\cup Z\cup \widehat{V}\Bigr)\;\#\;
\Bigl(E_{J}\cup V\Bigr),
\]
where:
\begin{itemize}
\item $V$ is the solid torus glued to $T=\partial E_{J}$ along the slope
represented on $T$ by the cabling class
$p[\lambda_{J}]+q[\mu_{J}]$; and
\item $\widehat{V}$ is the solid torus glued to the $T$-boundary of
$\mathcal{C}\cup Z$ along the same slope, so that
$\mathcal{C}\cup Z\cup\widehat{V}$ is the ``pattern piece.''
\end{itemize}
By Convention~\ref{conv:surgery}, the slope
$p[\lambda_{J}]+q[\mu_{J}]$ on $\partial E_{J}$ is exactly the surgery
coefficient $q/p$; hence
\[
E_{J}\cup V=S^{3}_{q/p}(J).
\]
For the pattern piece, one observes that
$\mathcal{C}\cup Z\cup\widehat{V}$ depends only on the pattern
$T(p,q)$ and the slope data, not on the knot $J$; taking $J=O$ in the
whole formula identifies the pattern piece as $S^{3}_{p/q}(O)=L(p,q)$ by
Moser's theorem, since $C_{p,q}(O)=T(p,q)$
(Lemma~\ref{lem:cable-basics}(ii)) and
$S^{3}_{q/p}(O)=L(q,p)$ (Lemma~\ref{lem:lens}).  This proves the
displayed formula.

Finally, we independently verify the slope bookkeeping, i.e.\ that the
companion coefficient must be $q/p$, by a first-homology computation
(which is a useful consistency check on all conventions).  By
Lemma~\ref{lem:H1},
\[
\bigl|H_{1}\bigl(S^{3}_{pq}(C_{p,q}(J))\bigr)\bigr|=|pq|,\qquad
\bigl|H_{1}\bigl(L(p,q)\bigr)\bigr|=p,\qquad
\bigl|H_{1}\bigl(S^{3}_{q/p}(J)\bigr)\bigr|=|q|,
\]
and the order of $H_{1}$ of a connected sum is the product of the orders
(Lemma~\ref{lem:H1}).  Thus the decomposition
$S^{3}_{pq}(C_{p,q}(J))\cong L(p,q)\#\,S^{3}_{q/p}(J)$ is consistent with
homology, and the companion coefficient is forced (given the lens-space
summand $L(p,q)$, whose order $p$ is the winding number) to have
numerator $\pm q$.  Together with the Moser anchor above and with the
classical winding law ``companion slope $=$ cabling slope $/$
$($winding$)^{2}$'' of the splicing calculus \cite{EN85}---here
$pq/p^{2}=q/p$---this pins down the formula as stated.
\end{proof}

\begin{remark}
Equivalently, Proposition~\ref{prop:cable-split} says that the companion
summand is the Dehn filling of $E_{J}$ along the cabling class
$p[\lambda_{J}]+q[\mu_{J}]$ itself.  Note that the coefficient $q/p$ is
\emph{never} an integer under the hypotheses $p\ge 2$,
$\gcd(p,|q|)=1$; this is exactly what will force the second summand to
be irreducible.
\end{remark}

We next recall the foundational theorem on non-integral surgeries.

\begin{theorem}[Culler--Gordon--Luecke--Shalen; see also
Gordon--Luecke]\label{thm:nonintegral}
Let $K\subset S^{3}$ be a nontrivial knot, and let $p/q$ be a slope in
lowest terms with $q\ge 2$.  Then $S^{3}_{p/q}(K)$ is irreducible.
Equivalently: if $S^{3}_{p/q}(K)$ is reducible and $K$ is nontrivial,
then $q=1$.
\end{theorem}

\begin{proof}
This is one of the foundational results of the theory of exceptional
surgeries, obtained in \cite{CGLS} from the analysis of essential planar
surfaces in knot exteriors; see also \cite{GL89}.  We use it as a black
box.
\end{proof}

\begin{lemma}\label{lem:q-over-p}
Let $p\ge2$ and $q\neq0$ with $\gcd(p,|q|)=1$.  Then $q/p$ is a
non-integral slope (in lowest terms its denominator is $p\ge2$).
\end{lemma}

\begin{proof}
If $q/p$ were an integer, then $p\mid q$, so $p\mid\gcd(p,|q|)=1$,
contradicting $p\ge2$.
\end{proof}

Finally, we record the standard consequence of the Gordon--Luecke
complement theorem used in the refinement of
Section~\ref{sec:remarks}.

\begin{theorem}[Gordon--Luecke {\cite{GL89}}]\label{thm:GL}
If $K\subset S^{3}$ is a nontrivial knot, then $S^{3}_{p/q}(K)\not\cong
S^{3}$ for every slope $p/q$; equivalently, if $S^{3}_{p/q}(K)\cong
S^{3}$ then $K=O$.
\end{theorem}

\section{Proof of the Main Theorem}\label{sec:proof}

We now assemble the ingredients.  For the reader's convenience we
restate the Main Theorem.

\begin{theorem}\label{thm:main}
Let $K\subset S^{3}$ be an $L$-space knot, and let $r\in
\mathbb{Q}\cup\{\infty\}$.  Then $S^{3}_{r}(K)$ has at most two prime
summands.  Moreover:
\begin{enumerate}
\item[(i)] if $S^{3}_{r}(K)$ is irreducible, it is prime;
\item[(ii)] if $S^{3}_{r}(K)$ is reducible, then $K=C_{p,q}(J)$ for some
knot $J\subset S^{3}$ and coprime $p\ge2$, $q\neq0$, one has $r=pq$, and
\[
S^{3}_{r}(K)\;\cong\;L(p,q)\;\#\;S^{3}_{q/p}(J),
\]
where $L(p,q)$ and $S^{3}_{q/p}(J)$ are prime.  Hence $S^{3}_{r}(K)$
has exactly two prime summands.
\end{enumerate}
\end{theorem}

\begin{proof}
If $r=\infty$ then $S^{3}_{r}(K)=S^{3}$, which has zero prime summands.
Henceforth $r\in\mathbb{Q}$, and we write
$Y:=S^{3}_{r}(K)$.

\medskip
\noindent\textbf{Step 0 (the unknot).}
Suppose first that $K=O$.  Then $Y$ is $S^{3}$, $S^{1}\times S^{2}$, or a
lens space (Lemma~\ref{lem:lens}); in each case $Y$ is prime or $S^{3}$,
so $Y$ has at most one prime summand.  Thus we may assume henceforth
that
\[
K\ \text{is a nontrivial } L\text{-space knot.}
\]

\medskip
\noindent\textbf{Step 1 (the irreducible case).}
If $Y$ is irreducible, then $Y$ is prime by
Lemma~\ref{lem:irr-prime}(i), so $Y$ has one prime summand (or zero, if
$Y\cong S^{3}$).  This proves (i).

\medskip
\noindent\textbf{Step 2 (the reducible case: cable structure).}
Suppose now that $Y$ is reducible, i.e.\ $Y\cong Y_{1}\#Y_{2}$ with
$Y_{i}\not\cong S^{3}$.  By Theorem~\ref{thm:HL} (Hedden--Levine), $K$
is a cable knot: there is a knot $J\subset S^{3}$ and coprime integers
$p\ge2$, $q\neq0$ with
\[
K=C_{p,q}(J)
\qquad\text{and}\qquad
r=pq .
\]
(This step is where the $L$-space hypothesis is used: for a general knot
the conclusion is precisely the open cabling conjecture.)

\medskip
\noindent\textbf{Step 3 (the splitting formula).}
By Proposition~\ref{prop:cable-split},
\[
Y=S^{3}_{pq}\bigl(C_{p,q}(J)\bigr)\;\cong\;
L(p,q)\;\#\;S^{3}_{q/p}(J).
\]
We claim that both summands are prime.

\emph{The first summand.}  $L(p,q)=S^{3}_{p/q}(O)$ is a lens space with
$|\pi_{1}L(p,q)|=p\ge2$, so $L(p,q)$ is irreducible
(Lemma~\ref{lem:lens}), hence prime
(Lemma~\ref{lem:irr-prime}(i)); in particular
$L(p,q)\not\cong S^{3}$.

\emph{The second summand.}  Let $M:=S^{3}_{q/p}(J)$.  There are two
cases.
\begin{itemize}
\item If $J=O$, then $M=L(q,p)$ is a lens space
(Lemma~\ref{lem:lens}); it is irreducible (hence prime), or $S^{3}$ if
$|q|=1$.  Note that if $J=O$ and $|q|=1$ then
$K=C_{p,\pm1}(O)=T(p,\pm1)$ is the unknot by
Lemma~\ref{lem:cable-basics}(ii), contradicting the nontriviality of
$K$; hence in fact $|q|\ge2$ in this case, and $M\cong L(q,p)\not\cong
S^{3}$.
\item If $J$ is nontrivial, then by Lemma~\ref{lem:q-over-p} the slope
$q/p$ is non-integral (its reduced denominator is $p\ge2$), so
Theorem~\ref{thm:nonintegral} implies that $M$ is irreducible.  By
Lemma~\ref{lem:irr-prime}(i), $M$ is prime, and by
Theorem~\ref{thm:GL}, $M\not\cong S^{3}$.
\end{itemize}
In both cases the second summand is prime and is not $S^{3}$ unless
possibly in the already-excluded degenerate situation.

\medskip
\noindent\textbf{Step 4 (conclusion).}
We have exhibited $Y$ as a connected sum of two prime manifolds, neither
of which is $S^{3}$.  By the uniqueness of the prime decomposition
\cite{Mi62,He76}, $Y$ has exactly two prime summands.  Together with
Step~1, this shows that $Y$ has at most two prime summands in all cases,
and proves (ii).
\end{proof}

\section{Corollaries, examples, and remarks}\label{sec:remarks}

\begin{corollary}[Exactly two, with identifications]\label{cor:exact}
Let $K$ be a nontrivial $L$-space knot and suppose $S^{3}_{r}(K)$ is
reducible.  Then $S^{3}_{r}(K)$ has exactly two prime summands, namely
the lens space $L(p,q)$ with
$\bigl|H_{1}(L(p,q))\bigr|=p$ equal to the winding number of the cable
$K=C_{p,q}(J)$, and the irreducible manifold $S^{3}_{q/p}(J)$ with
$\bigl|H_{1}(S^{3}_{q/p}(J))\bigr|=|q|$.
\end{corollary}

\begin{proof}
Immediate from Steps 3--4 of the proof of Theorem~\ref{thm:main},
including the verification that neither summand is $S^{3}$: for $J$
nontrivial this is Theorem~\ref{thm:GL}, and for $J=O$ we have
$|q|\ge2$ since $K$ is nontrivial
(Lemma~\ref{lem:cable-basics}(ii)), so $S^{3}_{q/p}(O)=L(q,p)\not\cong
S^{3}$.
\end{proof}

\begin{corollary}[Hyperbolic $L$-space knots]\label{cor:hyperbolic}
Let $K$ be a hyperbolic $L$-space knot (e.g.\ the $(-2,3,7)$ pretzel
knot).  Then every Dehn surgery on $K$ is irreducible; in particular
every such surgery is prime.
\end{corollary}

\begin{proof}
If some surgery on $K$ were reducible, then Theorem~\ref{thm:HL} would
exhibit $K$ as a cable knot, contradicting hyperbolicity (a cable is a
satellite knot).  Hence no surgery is reducible, and
Theorem~\ref{thm:main}(i) applies to all slopes.
\end{proof}

\begin{example}\label{ex:cable-of-trefoil}
Let $J=T(2,3)$ be the right-handed trefoil, an $L$-space knot of genus
$1$, and let $K=C_{5,2}(J)$, the $(5,2)$-cable; since $p/q=5/2>
2g(J)-1=1$, Hedden's cabling theorem shows that $K$ is an $L$-space
knot.  By Theorem~\ref{thm:main}, the unique reducible surgery on $K$
is the cabling-slope surgery, and
\[
S^{3}_{10}(K)\;\cong\;L(5,2)\;\#\;S^{3}_{2/5}\bigl(T(2,3)\bigr).
\]
The first summand is a lens space; the second is irreducible by
Theorem~\ref{thm:nonintegral} because $2/5$ is non-integral and $T(2,3)$
is nontrivial.  Thus $S^{3}_{10}(K)$ has exactly two prime summands, while
every other surgery on $K$ is irreducible.
\end{example}

\begin{remark}[Relation to the cabling conjecture]
For an arbitrary knot $K\subset S^{3}$, the statement ``every surgery on
$K$ has at most two prime summands'' is, in view of the classical
calculus, essentially equivalent to the cabling conjecture: the
conjecture asserts that a reducible surgery forces $K$ to be a cable
and the slope to be the cabling slope, in which case
Proposition~\ref{prop:cable-split} produces exactly two prime summands
(as in Corollary~\ref{cor:exact}).  Theorem~\ref{thm:main} establishes
this conclusion unconditionally for the class of $L$-space knots, the
crucial input being Theorem~\ref{thm:HL}.
\end{remark}

\begin{remark}[Why two rather than three]
It is worth isolating why a \emph{third} prime summand can never appear
in the reducible case, even for iterated cables.  By
Proposition~\ref{prop:cable-split}, the second summand
$S^{3}_{q/p}(J)$ is a surgery on the companion $J$ at the slope
$q/p$, and by Lemma~\ref{lem:q-over-p} this slope is \emph{never
integral}, since $p\ge2$ and $\gcd(p,|q|)=1$.  Consequently:
\begin{itemize}
\item if $J$ is nontrivial, Theorem~\ref{thm:nonintegral} forces
$S^{3}_{q/p}(J)$ to be irreducible, so it contributes exactly one prime
summand, no matter how complicated $J$ is (even if $J$ is itself an
iterated cable with its own reducible integral surgeries---those occur at
integral slopes, never at $q/p$);
\item if $J=O$, the summand is the lens space $L(q,p)$.
\end{itemize}
Thus the naive possibility of ``chaining'' reducible surgeries along the
JSJ tower of an iterated cable is blocked, at every stage, by the
non-integrality of the induced companion slope.
\end{remark}

\begin{remark}[Counting convention]
Our convention that $S^{3}$ has zero prime summands and that
$S^{1}\times S^{2}$ is prime (Lemma~\ref{lem:irr-prime}(ii)) is the
standard one; with it, the statement ``at most two prime summands''
covers all closed manifolds obtained by Dehn surgery, including the
degenerate cases.  Note that $S^{1}\times S^{2}$ can never occur as one of
the two summands in the reducible case of
Theorem~\ref{thm:main}(ii), since $|H_{1}|=|pq|<\infty$ and
$H_{1}(S^{1}\times S^{2})\cong\mathbb{Z}$.
\end{remark}

\section*{Conclusion}

The answer to the question is \textbf{yes}: a $3$-manifold obtained by
Dehn surgery on an $L$-space knot has at most two prime summands.  If
the surgery is reducible, the knot is necessarily a cable $C_{p,q}(J)$,
the slope is the cabling slope $pq$, and the output is the connected sum
of the lens space $L(p,q)$ with the irreducible manifold
$S^{3}_{q/p}(J)$; otherwise the output is irreducible, hence prime.

\begin{thebibliography}{10}

\bibitem{CGLS}
M. Culler, C.~M. Gordon, J. Luecke, and P.~B. Shalen,
\emph{Dehn surgery on knots},
Invent. Math. \textbf{87} (1987), 1--51.

\bibitem{EN85}
D. Eisenbud and W. Neumann,
\emph{Three-dimensional link theory and invariants of plane curve
singularities},
Annals of Mathematics Studies \textbf{110}, Princeton University Press,
Princeton, NJ, 1985.

\bibitem{GL89}
C.~M. Gordon and J. Luecke,
\emph{Knots are determined by their complements},
J. reine angew. Math. \textbf{399} (1989), 1--39.

\bibitem{He76}
J. Hempel,
\emph{$3$-manifolds},
Annals of Mathematics Studies \textbf{86}, Princeton University Press,
Princeton, NJ, 1976.

\bibitem{HL16}
M. Hedden and A.~S. Levine,
\emph{Splicing knot complements and bordered Heegaard Floer homology},
J. reine angew. Math. \textbf{720} (2016), 129--159.

\bibitem{Mi62}
J. Milnor,
\emph{A unique decomposition theorem for $3$-manifolds},
Amer. J. Math. \textbf{84} (1962), 1--7.

\bibitem{Mos71}
L. Moser,
\emph{Elementary surgery along a torus knot},
Pacific J. Math. \textbf{38} (1971), 537--545.

\bibitem{OSz04a}
P. Ozsv\'ath and Z. Szab\'o,
\emph{Holomorphic disks and topological invariants for closed
three-manifolds},
Ann. of Math. (2) \textbf{159} (2004), 1027--1158.

\bibitem{OSz04b}
P. Ozsv\'ath and Z. Szab\'o,
\emph{Holomorphic disks and knot invariants},
Adv. Math. \textbf{186} (2004), 58--116.

\bibitem{OSz05}
P. Ozsv\'ath and Z. Szab\'o,
\emph{On knot Floer homology and lens space surgeries},
Topology \textbf{44} (2005), 1281--1300.

\bibitem{Ro90}
D. Rolfsen,
\emph{Knots and links},
Corrected reprint of the 1976 original, Mathematics Lecture Series
\textbf{1}, Publish or Perish, Houston, TX, 1990.

\end{thebibliography}

\end{document}
```

**Summary of the argument.**  The proof rests on two cited pillars and otherwise proceeds by complete elementary arguments.  Pillar 1 is the Hedden--Levine theorem (the cabling conjecture for $L$-space knots): a reducible surgery on a nontrivial $L$-space knot forces the knot to be a cable $C_{p,q}(J)$ and the slope to be the cabling slope $pq$.  Pillar 2 is the classical cable surgery calculus (anchored by Moser's theorem for torus knots), which gives the splitting
$S^{3}_{pq}(C_{p,q}(J)) \cong L(p,q)\,\#\,S^{3}_{q/p}(J)$, with the companion coefficient $q/p$ pinned down by the first-homology factorization $pq = p\cdot q$ and the Moser anchor.  The conclusion then follows from two fully proved observations: a lens space is prime, and the slope $q/p$ is never an integer (since $p \ge 2$, $\gcd(p,q)=1$), so the second summand is irreducible---either by the CGLS non-integrability theorem when $J$ is nontrivial, or because it is a lens space when $J$ is the unknot.  Hence every surgery on an $L$-space knot has at most two prime summands, with exactly two (a lens space and an irreducible manifold) precisely in the reducible case.